Spoľahlivá lokalizácia predstavuje jeden zo základných predpokladov pre úspešnú autonómnu navigáciu mobilných robotov v~reálnom prostredí. Vzhľadom na obmedzenia a~neistoty jednotlivých snímačov, ako aj potenciálne výpadky v~reálnych aplikáciách, sa fúzia údajov zo snímačov javí ako kľúčový nástroj na zvýšenie presnosti a~robustnosti odhadu polohy. Táto práca poskytuje systematický prehľad dostupných senzorických technológií, ktoré sa využívajú na určenie polohy mobilných robotov a~analyzuje ich z~hľadiska princípov fungovania, vlastností a~možností aplikácie. Opísané sú tiež zaužívané pravdepodobnostné metódy senzorickej fúzie, ako aj súčasné trendy smerujúce k~využitiu techník umelej inteligencie pri spracovaní a~integrácii senzorických údajov. Výsledky analytickej časti práce tvoria teoretický základ pre návrh nových postupov zvyšujúcich spoľahlivosť lokalizačných systémov bez potreby externej kalibračnej infraštruktúry, čím prispievajú k~ďalšiemu rozvoju autonómnych robotických technológií v~širokej škále odvetví vrátane priemyslu či logistiky.